%% =============================================================================
%% DEFINING CHAT-FIRST ANALYTICS
%% =============================================================================

\section{Background and design principles} \label{sec:defining}

The introduction identified three concerns with using language models for
data analysis: correctness, privacy, and reproducibility. This section
examines these failure modes more carefully, then states the three
principles that define chat-first analytics.

\subsection{How LLM-assisted analysis goes wrong}

The appeal of typing ``regress wages on education with robust standard
errors'' into ChatGPT is obvious. The problems are less visible. They
fall into four categories.

First, LLM-generated code often runs without error but computes the wrong
quantity. Requesting ``clustered standard errors'' might yield
heteroskedasticity-robust errors instead. A panel model might omit
entity fixed effects. A difference-in-differences specification might
use the wrong interaction term. \citet{Ludwig+Mullainathan+Rambachan:2025}
document these risks systematically: the code executes, returns
plausible-looking output, and the user has no way to detect the
error without the expertise that would have let them write the code
themselves.

Second, code-generation systems require the data. When a user asks ChatGPT
to analyze a dataset, the data are uploaded to a third-party server
with unclear retention policies. For survey responses, health records,
personnel data, or any dataset subject to privacy regulations, this
may be unacceptable. Even when providers promise not to retain data,
the transmission itself may violate institutional or regulatory
requirements.

Third, the same prompt sent to the same model on different days may produce
different code. Models are updated, fine-tuned, and retrained. The
generated \proglang{Python} script that worked in January may not be
what the model produces in March. Conversational sessions are
ephemeral---unless the user manually saves the generated code, the
analysis cannot be reproduced. This contradicts a basic requirement
of scientific work.

Fourth, when a system generates 30 lines of \proglang{Python} to answer a
statistical question, understanding what was computed requires reading
and auditing those 30 lines of code. The user who could not write the code
cannot audit it either. This problem extends beyond research:
enterprise text-to-SQL deployments struggle with the same issue---ambiguous
questions admit multiple valid interpretations, and users cannot verify
which was executed~\citep{Renggli+Ilyas+Rekatsinas:2025}.

These four problems are distinct but related. Silent errors arise
because the language model controls both \emph{what} to compute and
\emph{how} to compute it. Privacy violations arise because
computation and data must be co-located on the provider's servers.
Reproducibility failures arise because the generated code is a
function of the model's current state. Auditability failures arise
because the computational logic is opaque---buried in generated code
rather than invoked from a known, inspectable implementation.
Systematic error rates for LLM-generated statistical code are not
yet well established: \citet{Ludwig+Mullainathan+Rambachan:2025}
document the \emph{types} of errors that arise, but base rates
depend on the model, the prompt, and the statistical method. The
motivation for the architecture described below rests on the
existence and nature of these failure modes, not on their precise
frequency.

\subsection{Three principles} \label{sec:principles}

Chat-first data analytics addresses these problems through three
principles. ``Chat-first'' means that natural language conversation
is the designed-for entry point: tool schemas are annotated for LLM
consumption, error messages are phrased for conversational context,
and default parameters are chosen so that typical requests need
minimal specification. It does not mean chat-only---\pkg{prompt2analytics}
also provides a CLI that invokes the same methods directly
(Section~\ref{sec:tools}), and every tool call maps to a public
\proglang{Rust} function usable without any LLM---but the
conversational interface is the primary design target. The first
two principles are definitional; the third is a design choice that we
consider essential.

\subsubsection{Principle 1: Separate interpretation from computation}

The language model decides \emph{what} to compute; validated software
decides \emph{how}. When a user asks to ``estimate the effect of
education on wages, controlling for experience, with robust standard
errors,'' the language model's job is to determine that this requires
OLS with HC1 standard errors, identify the column names, and generate
a structured tool call. The actual matrix algebra---forming
$(\bm{X}'\bm{X})^{-1}\bm{X}'\bm{y}$, computing the
HC1 sandwich estimator---is performed by pre-implemented code that
the language model cannot modify.

This division of labor plays to each component's strengths. Language
models are good at understanding intent, resolving ambiguity, and
mapping natural language to structured specifications. They are
unreliable at generating correct numerical code. Separation removes
the failure mode where the model produces code that runs but computes
the wrong quantity, because the model never produces code at all.
Wrong parameter extraction remains possible (Section~\ref{sec:evaluation}),
but the error surface is smaller: the model selects from a fixed schema
rather than generating arbitrary code.

The idea has precedent in text-to-SQL systems~\citep{Yu+etal:2018,
Li+etal:2024}, where a language model translates natural language into
structured database queries executed by a SQL engine. We extend this
pattern from data retrieval to statistical computation: from
text-to-SQL to text-to-statistics.

\subsubsection{Principle 2: Pre-validate methods, not generated code}

Correctness is established once, through systematic validation against
reference implementations, rather than hoped for anew with each
request. When \pkg{prompt2analytics} runs OLS, it uses the same
\proglang{Rust} function every time---a function validated against
\proglang{R}'s \fct{lm} on the Longley dataset to $10^{-8}$
precision (Section~\ref{sec:evaluation}). The user need not audit
the implementation on each invocation, any more than they audit
\proglang{R}'s source code when calling \fct{lm}.

This principle also solves the reproducibility problem. The same
method invocation with the same data produces identical results,
regardless of when it is run or which language model is used. The
stochastic element of the language model affects only interpretation
(which tool to call, how to explain the output), never computation.
Two users who phrase the same analysis differently may get different
tool calls---but once the tool is selected and parameters specified,
the result is deterministic.

\subsubsection{Principle 3: Keep data local; transmit only structure}

The language model receives column names, data types, and summary
statistics---not the data itself. A dataset of 100,000 salary records
is represented to the model as metadata: ``columns: employee\_id (int),
salary (float, range 28000--185000), department (string, 12 unique
values), years\_experience (float, range 0.5--35.2).'' Statistical
computations execute on the user's machine; only aggregated results
(coefficients, standard errors, test statistics) pass through the
language model.

This is not definitionally required for chat-first analytics---a
cloud-based system satisfying Principles~1 and~2 would still qualify.
But we consider data locality essential for the use cases that
motivate this work. Researchers analyzing survey data, firms analyzing
personnel records, and analysts working under privacy regulations need
architectures where data isolation is guaranteed by design, not
promised by policy.

For complete isolation, the language model itself can run locally.
Open-source models deployed via Ollama~\citep{Ollama} eliminate all
external network communication. Our evaluation
(Section~\ref{sec:evaluation}) shows that even a 3-billion-parameter
model achieves over 80\% adequate rate, making fully local
deployment practical on consumer hardware.

Together, the three principles address the four failure modes
identified above. Separation and pre-validation eliminate silent
computational errors: the model never generates code, and the code
it invokes is validated once against reference implementations.
Pre-validation also ensures reproducibility, because the same method
invocation with the same data always produces the same result. Data
locality protects privacy by keeping raw data on the user's machine.
Auditability improves on all three fronts---the user inspects a
structured tool call rather than 30 lines of generated code.

\subsection{Relationship to existing ideas}

Chat-first analytics applies established techniques to a new domain.
Tool-augmented language models~\citep{Schick+etal:2023, Yao+etal:2023}
can reliably select and invoke external tools when provided with
appropriate schemas~\citep{Qu+etal:2025, Liu+etal:2025:ToolACE}. The
Model Context Protocol (MCP)~\citep{MCP:2024} standardizes how
language models discover and invoke tools, providing the communication
layer between interpretation and computation.%
\footnote{MCP focuses on model-to-tool integration. Complementary
protocols address agent-to-agent communication: Google's Agent2Agent
(A2A) protocol~\citep{A2A:2025} and IBM's Agent Communication
Protocol (ACP)~\citep{ACP:2025}, which merged under Linux Foundation
governance in 2025. For single-agent tool invocation---the pattern
used here---MCP is the dominant standard, with adoption by OpenAI,
Microsoft, and Google alongside Anthropic.}

The interface innovation has precedents in four strands of work.
First, \emph{formula notation and domain-specific languages}:
the Wilkinson-Rogers formula notation~\citep{Wilkinson+Rogers:1973}
simplified model specification in the 1970s;
\citet{Chambers+Hastie:1992} extended it in \proglang{S} and
\proglang{R}~\citep{Chambers:1998}; and probabilistic programming
languages like Stan~\citep{Carpenter+etal:2017} apply the DSL
pattern to Bayesian inference. Chat-first analytics replaces
formal notation with natural language, trading precision for
accessibility.
Second, \emph{natural language interfaces to data systems}:
LUNAR~\citep{Woods:1973} demonstrated feasibility for database
querying in the 1970s, and the text-to-SQL paradigm---driven by
benchmarks like Spider~\citep{Yu+etal:2018}---has produced
substantial improvements in translating natural language to
structured queries~\citep{Li+etal:2024}. We extend this pattern
from data retrieval to statistical computation.
Third, \emph{automated machine learning}: Auto-sklearn~\citep{Feurer+etal:2015}
and H2O AutoML~\citep{LeDell+Poirier:2020} automate model selection
through algorithmic search, and recent work integrates LLMs into
the pipeline~\citep{Hollmann+etal:2024}. AutoML optimizes for
prediction; chat-first analytics interprets user intent across
prediction, causal inference, hypothesis testing, and description.
Fourth, \emph{computational econometrics} has established
standards for numerical accuracy~\citep{McCullough:1997,
McCullough+Vinod:1999, Davidson+MacKinnon:2004}, which motivate
the validation approach described in Section~\ref{sec:evaluation}.

Several systems have explored LLM-assisted data analysis, but with
a code-generation architecture. \citet{Hong+etal:2024} present a
``Data Interpreter'' that generates and executes \proglang{Python}
code. \citet{Chen+Shen+Ma:2025} develop an ``Econometrics AI Agent''
that augments code generation with domain knowledge.
\citet{Sun+etal:2025:LAMBDA} describe an agent that orchestrates
generated code across analysis stages. All three generate code;
none pre-validate implementations or enforce data locality.
Table~\ref{tab:system-comparison} summarizes these differences.

Hybrid designs are possible. A system
could generate code for bespoke analyses while routing standard
requests to pre-validated tools. It could verify generated code
through static analysis or property checking before execution---for
instance, confirming that a difference-in-differences specification
includes the correct interaction term, or that a panel regression
includes entity fixed effects when requested. It could also combine
code generation with human review, leveraging the fact that many users
can \emph{audit} code even if they cannot \emph{write} it from
scratch. \pkg{prompt2analytics} deliberately occupies one end of
this spectrum: all computation is pre-validated, with no code
generation at all. This sacrifices flexibility---analyses outside the
method library are unavailable, custom data transformations beyond the
built-in munging tools require the CLI or direct \proglang{Rust}
code, and ad-hoc visualizations not covered by the visualization
tools cannot be produced---but it eliminates an entire class of
failure modes. Whether hybrid designs can achieve a better trade-off
is an open question; the answer likely depends on the reliability of
code verification techniques and the target user's ability to review
generated code.

\begin{table}[htbp]
\centering
\begin{threeparttable}
\caption{Comparison with code-generation approaches}
\label{tab:system-comparison}
\small
\begin{tabular}{@{}lp{2.5cm}p{2.4cm}p{2.4cm}p{2.4cm}@{}}
\toprule
 & \pkg{p2a} & Econ.\ AI & Data Interpr. & LAMBDA \\
\midrule
LLM role & Tool selection & Enhanced code gen & Code gen + exec & Agent orchestration \\
Computation & Pre-validated & LLM-generated & LLM-generated & LLM-generated \\
Local execution & Yes (Ollama) & No & No & No \\
Validated & Yes (vs.\ R) & No & No & No \\
Method flexibility & Fixed library (${\sim}200$ methods) & Any \proglang{R} method & Any \proglang{Python} method & Any \proglang{Python} method \\
Extensibility & Requires Rust impl. & Extends with model & Extends with model & Extends with model \\
\bottomrule
\end{tabular}
\begin{tablenotes}[flushleft]
\item \emph{Note:} Econ.\ AI = Econometrics AI
Agent~\citep{Chen+Shen+Ma:2025}; Data Interpr.\ = Data
Interpreter~\citep{Hong+etal:2024}; LAMBDA~\citep{Sun+etal:2025:LAMBDA}.
The trade-off: code-generation systems offer unlimited method
flexibility at the cost of unvalidated computation; \pkg{p2a}
constrains coverage to ensure correctness.
\end{tablenotes}
\end{threeparttable}
\end{table}

Three limitations of the approach are fundamental, not fixable by
better engineering. First, statistical expertise is still required:
chat-first analytics assists with method selection and execution, not
with research design or causal identification. A user who requests
``estimate the effect of X on Y'' receives a regression, but the
system does not validate whether a causal interpretation is warranted.
LLM explanations may echo causal language from the user's prompt even
when the analysis provides only associational evidence. Users remain
responsible for ensuring that their identification strategy supports
their claims.

Second, method selection is imperfect. Tool selection works well on
curated prompts (Section~\ref{sec:evaluation}), but users should
verify that the system understood their request correctly---particularly
for ambiguous specifications where multiple methods could apply. Users
should consult the accuracy results before interpreting the curated
examples in Section~\ref{sec:examples}.

Third, pre-validation trades coverage for correctness. When a user
requests a method outside the implemented set---Bayesian inference,
structural equation modeling, or a custom estimator---the LLM may
silently map the request to the closest available alternative rather
than reporting unavailability---the symmetric failure mode to
code generation's silent errors: instead of running wrong code, the
system runs the wrong method. The risk is mitigated by the structured
tool schema (which makes the selected method visible to the user) and
by the LLM's ability to explain substitutions, but users must still
verify that the method matches their intent.

These limitations mean that chat-first analytics is a useful tool,
not an autonomous analyst. Users should verify tool selection, check
extracted parameters, and apply domain knowledge to interpret
results---the same discipline required when using any statistical
software, with the added need to verify that the system understood the
request correctly.
